% =============================================================================
% --- secciones/4-ModeloDOT.tex ---
% =============================================================================

\section{Modelo DOT}

Frente al colapso del modelo de expansión dispersa, el 
\textbf{Desarrollo Orientado al Transporte (DOT)} surge como la estrategia de regeneración 
urbana más robusta para las metrópolis mexicanas. 
Su objetivo central es revertir la dependencia del automóvil mediante la creación de núcleos 
de actividad compactos, de usos mixtos y alta densidad, articulados alrededor de estaciones 
de transporte público masivo.

\subsection{Los Ocho Principios del Diseño Urbano Sustentable}
\label{subsubsec:ocho-principios}

El ITDP ha sistematizado la implementación del modelo DOT a través de ocho principios 
rectores que sirven como marco de evaluación para el urbanista. Estos principios no solo 
buscan la eficiencia técnica, sino la vitalidad del tejido social:

\begin{itemize}
    \item \textbf{Caminar y Pedalear:} Establecer la movilidad activa como la base de 
            toda intermodalidad, garantizando entornos seguros y accesibles.
    \item \textbf{Conectar:} Crear redes de calles densas que permitan trayectos directos 
            para peatones y ciclistas.
    \item \textbf{Transportar, Densificar y Compactar:} Maximizar la densidad habitacional 
            y de empleo en las zonas con mejor servicio de transporte masivo, 
            limitando el crecimiento fuera de la huella urbana existente.
    \item \textbf{Mezclar y Cambiar:} Fomentar la mixticidad de usos de suelo y reducir 
            activamente el espacio dedicado al vehículo particular.
\end{itemize}

Para transitar de la descripción cualitativa a la evaluación rigurosa, el \textbf{Estándar DOT} 
utiliza un sistema de puntuación que reconoce proyectos con niveles de 
\textbf{Oro, Plata o Bronce} \cite{batista2015, ugur2025}. 
Esta métrica es fundamental para garantizar que la densificación no resulte en procesos de 
gentrificación, sino en una auténtica \textbf{justicia socioespacial} que preserve 
la vivienda asequible cerca de la infraestructura masiva \cite{moreno2023}.

\subsection{Reporte 02: La Brecha de Cobertura Capilar ($C$)}
\label{subsubsec:reporte-02}

La viabilidad del modelo DOT en la ZMVM está condicionada por la equidad en la distribución 
de la red masiva. El \textbf{Modelo VFT} crea una operación para esta condición mediante 
el indicador de \textbf{Nivel de Cobertura ($C$)}, el cual evalúa la intersección 
entre el área de influencia de las estaciones ($A_{cubierta}$) y la superficie total 
de la demarcación ($A_{total}$) \cite{Garibaldi2024VFT, Garibaldi2024_R02}:

\begin{equation}
    C = \left( \frac{A_{cubierta}}{A_{total}} \right) \times 100
\end{equation}

Siguiendo los lineamientos del ITDP, se estableció un 
\textbf{isócrono de accesibilidad peatonal de 800 metros} 
(equivalente a 10-15 minutos de caminata) como el umbral de 
cobertura efectiva. 

\subsection{Contraste de Realidad: El Agotamiento del Modelo Radial}
\label{subsubsec:agotamiento-radial}

El análisis comparativo por alcaldía revela que el modelo radial ha alcanzado 
un estado de saturación en el centro, mientras que las periferias sufren una 
exclusión sistemática. Mientras que el \termino{Real Case} 
de la alcaldía \textbf{Benito Juárez} presenta niveles de cobertura capilar 
superiores al \textbf{79\%-98\%} debido a la redundancia de sistemas masivos 
(Metro, Metrobús, Trolebús), las zonas del tercer y cuarto contorno muestran un 
rezago crítico.

Los datos del Reporte 02 confirman que en demarcaciones como \textbf{Milpa Alta (11.75\%)} 
y \textbf{Tlalpan (30.01\%)} el principio de accesibilidad universal es inexistente. 
Este hallazgo demuestra que el modelo DOT no puede ser implementado con éxito de forma 
aislada; requiere de una \textbf{interceptación transversal} que rompa la inercia del 
centro saturado. La propuesta del \textbf{Corredor Anillar} se justifica técnicamente 
como el mecanismo para cerrar esta brecha de cobertura, permitiendo que los beneficios 
del DOT se extiendan hacia los territorios históricamente segregados.
